\documentclass[10pt,a4paper]{amsart}
\usepackage{amsmath,amssymb,amsthm,graphicx}
% Retained preamble for the Stacks Project chapter snapshot in this directory.
% The chapter's own preamble.tex is a copy of the shared Stacks Project preamble; the block
% below is sliced MECHANICALLY, line for line, from it. Nothing was added, removed or
% redefined; the wrapper class, the xr-hyper/hyperref wiring and the externaldocument list
% are omitted because the assembler supplies the document class.
%
% lines taken from preamble.tex: 3, 5, 9, 11, 12, 156, 157, 158, ... (51 lines total)
%
% package coverage restored from the chapter's own preamble (2026-10-13): xy, xy-2cell

\usepackage[all]{xy}
\xyoption{2cell}
\UseAllTwocells

\newenvironment{reference}{\comment}{\endcomment}
\newenvironment{slogan}{\comment}{\endcomment}
\newenvironment{history}{\comment}{\endcomment}
\theoremstyle{plain}
\newtheorem{theorem}[subsection]{Theorem}
\newtheorem{proposition}[subsection]{Proposition}
\newtheorem{lemma}[subsection]{Lemma}
\theoremstyle{definition}
\newtheorem{definition}[subsection]{Definition}
\newtheorem{example}[subsection]{Example}
\newtheorem{exercise}[subsection]{Exercise}
\newtheorem{situation}[subsection]{Situation}
\theoremstyle{remark}
\newtheorem{remark}[subsection]{Remark}
\newtheorem{remarks}[subsection]{Remarks}
\numberwithin{equation}{subsection}
\def\lim{\mathop{\mathrm{lim}}\nolimits}
\def\colim{\mathop{\mathrm{colim}}\nolimits}
\def\Spec{\mathop{\mathrm{Spec}}}
\def\Hom{\mathop{\mathrm{Hom}}\nolimits}
\def\Ext{\mathop{\mathrm{Ext}}\nolimits}
\def\SheafHom{\mathop{\mathcal{H}\!\mathit{om}}\nolimits}
\def\SheafExt{\mathop{\mathcal{E}\!\mathit{xt}}\nolimits}
\def\Sch{\mathit{Sch}}
\def\Mor{\mathop{\mathrm{Mor}}\nolimits}
\def\Ob{\mathop{\mathrm{Ob}}\nolimits}
\def\Sh{\mathop{\mathit{Sh}}\nolimits}
\def\NL{\mathop{N\!L}\nolimits}
\def\CH{\mathop{\mathrm{CH}}\nolimits}
\def\proetale{{pro\text{-}\acute{e}tale}}
\def\etale{{\acute{e}tale}}
\def\QCoh{\mathit{QCoh}}
\def\Ker{\mathop{\mathrm{Ker}}}
\def\Im{\mathop{\mathrm{Im}}}
\def\Coker{\mathop{\mathrm{Coker}}}
\def\Coim{\mathop{\mathrm{Coim}}}
\DeclareMathSymbol{\boxtimes}{\mathbin}{AMSa}{"02}
\def\QCohstack{\mathcal{QC}\!\mathit{oh}}
\def\Cohstack{\mathcal{C}\!\mathit{oh}}
\def\Spacesstack{\mathcal{S}\!\mathit{paces}}
\def\Quotfunctor{\mathrm{Quot}}
\def\Hilbfunctor{\mathrm{Hilb}}
\def\Curvesstack{\mathcal{C}\!\mathit{urves}}
\def\Polarizedstack{\mathcal{P}\!\mathit{olarized}}
\def\Complexesstack{\mathcal{C}\!\mathit{omplexes}}
\def\Pic{\mathop{\mathrm{Pic}}\nolimits}
\def\Picardstack{\mathcal{P}\!\mathit{ic}}
\def\Picardfunctor{\mathrm{Pic}}
\def\Deformationcategory{\mathcal{D}\!\mathit{ef}}
\usepackage{cleveref}
\crefname{originalequation}{original equation}{original equations}
\Crefname{originalequation}{Original equation}{Original equations}
\makeatletter
\newlabel{equation-fibre-product}{{09DK}{source}{}{sourcepointer-1}{}}
\newlabel{equation-fibre-product@cref}{{[originalequation][0][]09DK}{[1][0][]source}{}{}{}}
\makeatother
\crefname{originallemma}{original lemma}{original lemmas}
\Crefname{originallemma}{Original lemma}{Original lemmas}
\makeatletter
\newlabel{lemma-compare-spaces-schemes}{{08VF}{source}{}{sourcepointer-2}{}}
\newlabel{lemma-compare-spaces-schemes@cref}{{[originallemma][0][]08VF}{[1][0][]source}{}{}{}}
\makeatother
\crefname{originallemma}{original lemma}{original lemmas}
\Crefname{originallemma}{Original lemma}{Original lemmas}
\makeatletter
\newlabel{lemma-morphism-affine-schemes}{{08T3}{source}{}{sourcepointer-3}{}}
\newlabel{lemma-morphism-affine-schemes@cref}{{[originallemma][0][]08T3}{[1][0][]source}{}{}{}}
\makeatother
\crefname{originallemma}{original lemma}{original lemmas}
\Crefname{originallemma}{Original lemma}{Original lemmas}
\makeatletter
\newlabel{lemma-tensor-product-tor-independent}{{09DA}{source}{}{sourcepointer-4}{}}
\newlabel{lemma-tensor-product-tor-independent@cref}{{[originallemma][0][]09DA}{[1][0][]source}{}{}{}}
\makeatother
\crefname{originallemma}{original lemma}{original lemmas}
\Crefname{originallemma}{Original lemma}{Original lemmas}
\makeatletter
\newlabel{spaces-perfect-lemma-tor-independent}{{08IQ}{source}{}{sourcepointer-5}{}}
\newlabel{spaces-perfect-lemma-tor-independent@cref}{{[originallemma][0][]08IQ}{[1][0][]source}{}{}{}}
\makeatother
\title{The cotangent complex of a Tor independent fibre product of algebraic spaces}
\author{The Stacks Project Authors}
\date{stacks-project:cotangent.tex:local-snapshot-20261012; GNU FDL 1.2 or later}
\begin{document}
\maketitle
\setcounter{subsection}{1}
\noindent\textit{Editorial note.} All delivered mathematical text is the author's own, copied from the pinned Stacks Project chapter file cotangent.tex; no statement or proof step was rewritten, reordered or removed.\par\medskip
\noindent\textit{Editorial note.} The statement refers to the situation and the fibre product equation fixed earlier in the chapter, and displays the resulting decomposition of the cotangent complex.\par\medskip
\noindent\textit{Editorial note.} The five references of the counted spans are recorded as non-load-bearing navigation pointers; the argument that consumes them stays inside the counted proof span.\par\medskip
\noindent\textit{Editorial note.} Preamble: the retained preamble of this chapter (sources/stacks-cotangent/retained-preamble.tex, lines 1-63), containing the Xy-pic setup, the comment environments, the theorem declarations and the macros including \textbackslash{}NL.\par\medskip
\noindent\textit{Allowed standard prerequisites (editorial).}\begin{itemize}
\item The comparison between the cotangent complexes of algebraic spaces and of their scheme presentations. Lemma lemma-compare-spaces-schemes, used at proof lines 4812-4814. Reduces the vanishing statement to schemes.
\item Tor independence for derived categories of algebraic spaces. Derived Categories of Spaces, Lemma spaces-perfect-lemma-tor-independent, used at proof lines 4812-4814. Transfers the Tor independence hypothesis to the presentation.
\item The comparison of the cotangent complex over an affine open with the ring-level complex. Lemma lemma-morphism-affine-schemes (4129-4140, produced as L44141), used at proof lines 4815-4819. Reduces the affine scheme case to a tensor product of rings.
\item The ring-level case. Lemma lemma-tensor-product-tor-independent, used at proof lines 4817-4819. Settles the final case.
\end{itemize}\medskip
\ifdefined\hypertarget\hypertarget{sourcepointer-1}{}\fi
\noindent\textit{Source reference.} Original equation 09DK: focused-seeds/AGP/sources/stacks-snapshot-20261011/tags.txt, source lines 11868--11868. Navigation pointer into the Stacks Project; the standard fact it supplies is applied inside the counted proof span and the argument resting on it stays there.\par\medskip
\ifdefined\hypertarget\hypertarget{sourcepointer-2}{}\fi
\noindent\textit{Source reference.} Original lemma 08VF: focused-seeds/AGP/sources/stacks-snapshot-20261011/tags.txt, source lines 11233--11233. Navigation pointer into the Stacks Project; the standard fact it supplies is applied inside the counted proof span and the argument resting on it stays there.\par\medskip
\ifdefined\hypertarget\hypertarget{sourcepointer-3}{}\fi
\noindent\textit{Source reference.} Original lemma 08T3: focused-seeds/AGP/sources/stacks-snapshot-20261011/tags.txt, source lines 11151--11151. Navigation pointer into the Stacks Project; the standard fact it supplies is applied inside the counted proof span and the argument resting on it stays there.\par\medskip
\ifdefined\hypertarget\hypertarget{sourcepointer-4}{}\fi
\noindent\textit{Source reference.} Original lemma 09DA: focused-seeds/AGP/sources/stacks-snapshot-20261011/tags.txt, source lines 11858--11858. Navigation pointer into the Stacks Project; the standard fact it supplies is applied inside the counted proof span and the argument resting on it stays there.\par\medskip
\ifdefined\hypertarget\hypertarget{sourcepointer-5}{}\fi
\noindent\textit{Source reference.} Original lemma 08IQ: focused-seeds/AGP/sources/stacks-snapshot-20261011/tags.txt, source lines 10823--10823. Navigation pointer into the Stacks Project; the standard fact it supplies is applied inside the counted proof span and the argument resting on it stays there.\par\medskip
\begin{lemma}
\label{lemma-fibre-product-tor-independent}
In the situation above, if $X$ and $Y$ are Tor independent over $B$, then
the object $E$ in (\ref{equation-fibre-product}) is zero. In this case we
have
$$
L_{X \times_B Y/B} = Lp^*L_{X/B} \oplus Lq^*L_{Y/B}
$$
\end{lemma}

\begin{proof}
Choose a scheme $W$ and a surjective \'etale morphism $W \to B$.
Choose a scheme $U$ and a surjective \'etale morphism $U \to X \times_B W$.
Choose a scheme $V$ and a surjective \'etale morphism $V \to Y \times_B W$.
Then $U \times_W V \to X \times_B Y$ is surjective \'etale too.
Hence it suffices to prove that the restriction of $E$ to $U \times_W V$
is zero. By Lemma \ref{lemma-compare-spaces-schemes} and
Derived Categories of Spaces, Lemma \ref{spaces-perfect-lemma-tor-independent}
this reduces us to the case of schemes.
Taking suitable affine opens we reduce to the case of affine schemes.
Using 
Lemma \ref{lemma-morphism-affine-schemes}
we reduce to the case of a tensor product of rings, i.e., to
Lemma \ref{lemma-tensor-product-tor-independent}.
\end{proof}

\end{document}
